% ------------------------------------------------------------------
% Section: No balanced rotor on nine-atom circuits (t=3)
% Drafted by: Claude (Fable 5), 2026-07-17
% Sources (all claims trace to these; see claude_t3_no_balanced_rotor_claims.md):
%   problems/23/writeup/WALL_ATTACK_R43_GPTPRO56.md         (secs. 1-2, 5-7)
%   problems/23/writeup/WALL_ATTACK_R42_GPTPRO56.md         (secs. 9-10, owner conditions)
%   problems/23/writeup/WALL_ATTACK_R41_GPTPRO56.md         (sec. 7, window parameters)
%   problems/23/writeup/WALL_ATTACK_R48_GPTPRO56.md         (secs. 1-4, profile conditions)
%   problems/23/archive/tmp_text_sources_20260712/fanout/
%     r42_graph_specific_exclusion/REPORT.md                (four-star lemma; bad-star
%       vertex-cover freeness; shape-independent t=3 closure; axiom audits)
%   problems/23/lean/Erdos23Delta0/Gamma/R43SupportIncidence.lean      (SHA-256 66649670...)
%   problems/23/lean/Erdos23Delta0/Gamma/BadStarCoverFreeness.lean     (SHA-256 afd944ea...)
%   problems/23/lean/Erdos23Delta0/Gamma/K33BadStarPairCountZero.lean  (SHA-256 35456523...)
%   problems/23/lean/Erdos23Delta0/Gamma/CutTightActiveRotorIncidence.lean (SHA-256 0814a665...)
% Cross-section dependencies (labels defined in sibling sections):
%   sec:rotor8, def:rotor8-selection   (claude_eight_vertex_rotor.tex)
%   cor:graph                          (support_circuits.tex)
%   def:switchloss, cor:badendpoint    (star_lemma.tex)
% Bibliography keys used: Ferudun26supports (as in support_circuits.tex).
% Assumes the companion preamble: amsthm environments, enumitem, \supp, \bip.
% ------------------------------------------------------------------

\section{No balanced rotor on nine-atom circuits
(\texorpdfstring{\(t=3\)}{t=3})}\label{sec:t3closure}

Section~\ref{sec:rotor8} exhibited a two-atom configuration, far from
any Hall violation, on which the local repair walk rotates forever.
This section treats the corresponding question on an actual support
circuit. The motivation is a defect-transport strategy for excluding
support circuits: one tries to show that some quantity strictly
decreases along detour moves, and the terminal objects such a strategy
must exclude are cyclic families of detour moves along which a locally
deficient vertex configuration merely rotates --- here with at most
three owners, the engine-relevant range on nine atoms
(Remark~\ref{rem:t3-fourowners}). We call these
\emph{balanced rotors} and define them precisely below. In the window
parametrization of the transport analysis, where a circuit of \(t^2\)
atoms carries rotating vertices with \(t\) crossing and \(t\)
monochromatic neighbours, the smallest case is \(t=3\): nine atoms and,
by the circuit identity, eight support edges. We give two independent
proofs that no balanced rotor exists on a nine-atom support circuit.
The first counts support edges and is insensitive to the shape of the
circuit; the second shows that on the unique nine-atom footprint the
coverage condition that any rotation needs is unsatisfiable. Both have
kernel-checked Lean counterparts (Remark~\ref{rem:t3-lean}).

Throughout the section, \(G\) is a finite triangle-free graph with a
fixed cut \(V(G)=V_0\mathbin{\dot\cup}V_1\), \(B\) is the crossing
subgraph, assumed connected, \(M=E(G)\setminus B\), and
\((F,\mathcal A)\) is a support circuit in the graph form of
Corollary~\ref{cor:graph}: \(\mathcal A\subseteq\{xy\in
M:d_B(x,y)=4\}\) is inclusion-minimal with
\(|\supp_B(\mathcal A)|<|\mathcal A|\), and
\(F=\supp_B(\mathcal A)\) is viewed as the bipartite subgraph it spans,
so that \(|E(F)|=|\mathcal A|-1\) and every \(B\)-geodesic of an atom
lies in \(F\). As in Definition~\ref{def:rotor8-selection}, a
\emph{selection} \(\omega\) assigns to each atom one of its
\(B\)-geodesics; \(S_\omega\subseteq E(F)\) denotes the union of the
selected edge sets, and a \emph{detour move} replaces the selected
geodesic of a single atom by a different geodesic of the same atom.
Maximality of the cut is not used in this section, except in the
motivational Remark~\ref{rem:t3-fourowners}.

\subsection{Rotating owners and balanced rotors}

\begin{definition}[Rotating owner profile; covered star]\label{def:t3-profile}
Let \(\omega\) be a selection. A vertex \(v\) has the \emph{rotating
owner profile} at \(\omega\) if:
\begin{enumerate}[label=\textup{(P\arabic*)}]
\item \(d_B(v)=d_M(v)=3\), and the three monochromatic edges at \(v\)
are atoms of \(\mathcal A\);
\item writing \(N_B(v)=\{x,y_1,y_2\}\), exactly one of the three
crossing edges at \(v\) lies outside \(S_\omega\); label the notation
so that this edge is \(vx\) (the \emph{active} edge), so
\(vy_1,vy_2\in S_\omega\);
\item (\emph{covered star}) for each \(i\in\{1,2\}\) some selected
geodesic contains both \(x\) and \(y_i\) but not \(v\).
\end{enumerate}
\end{definition}

\begin{remark}\label{rem:t3-profile-scope}
The definition keeps only the conditions used in the proofs below. In
the transport engine, a rotor state satisfies strictly stronger local
conditions (prescribed counts of selected geodesics through the owner
and two exact matching conditions), and its transitions carry an
exactly balanced deficiency ledger, which is the origin of the name
\emph{balanced} in Definition~\ref{def:t3-rotor}. Excluding the weaker
object above therefore suffices. Requiring in (P3) that the covering
geodesic avoid \(v\) also costs nothing: if a selected geodesic
\(w_0w_1w_2w_3w_4\) contained both \(v\) and \(x\), then, positions of
\(v\) and \(x\) having opposite parity because \(vx\in B\), either they
are adjacent on the geodesic, making \(vx\) a selected edge and
contradicting (P2), or they sit at positions \(\{0,3\}\) or
\(\{1,4\}\), and the chord \(vx\) yields a \(B\)-path of length two
between the atom endpoints, contradicting \(d_B(w_0,w_4)=4\).
\end{remark}

\begin{definition}[Balanced rotor]\label{def:t3-rotor}
A \emph{balanced rotor} on \((F,\mathcal A)\) consists of
\(k\in\{2,3\}\) distinct vertices \(v_1,\dots,v_k\) (the
\emph{rotating owners}) together with selections
\(\omega_1,\dots,\omega_k\) such that, with indices modulo \(k\):
\begin{enumerate}[label=\textup{(R\arabic*)}]
\item \(v_i\) has the rotating owner profile at \(\omega_i\);
\item \(\omega_{i+1}\) is obtained from \(\omega_i\) by a single detour
move replacing a selected geodesic of the form
\(a\,x\,v_{i+1}\,y\,b\) by the geodesic \(a\,x\,v_i\,y\,b\) of the same
atom \(ab\in\mathcal A\).
\end{enumerate}
\end{definition}

Condition (R2) says that consecutive owners occur as the two middle
vertices of a pair of geodesics of one atom that agree elsewhere; this
is exactly the one-atom detour move of the repair walk, as on the
four-state rotor of Section~\ref{sec:rotor8}.

\begin{remark}[Why two or three owners]\label{rem:t3-fourowners}
A nontrivial rotation needs at least two owners, and the engine's
four-owner configuration is excluded on nine atoms by a separate count,
which we record for orientation; it is the only place in this section
where maximality of the cut enters. In that configuration the four
rotating vertices are the corners of a crossing four-cycle, each is
cut-tight in the sense of Definition~\ref{def:switchloss}, and each has
at least four crossing neighbours. The four monochromatic stars are
pairwise disjoint: adjacent corners are joined by crossing edges, while
opposite corners lie on one shore and have common crossing neighbours,
so an edge between them would close a triangle.
Corollary~\ref{cor:badendpoint} gives \(d_M\ge d_B-1\ge3\) at each
corner, whence \(|M|\ge12>9\). The finite core of this count is
kernel-checked (Remark~\ref{rem:t3-lean}).
Definition~\ref{def:t3-rotor} therefore fixes \(k\in\{2,3\}\).
\end{remark}

\subsection{First closure: the support-edge count}

\begin{lemma}[Owner degree in the support graph]\label{lem:t3-ownerdeg}
Let \(\omega\) be a selection and let \(v\) have the rotating owner
profile at \(\omega\). Then all three crossing edges
\(vx,vy_1,vy_2\) are edges of \(F\); in particular \(v\in V(F)\) and
\(\deg_F(v)\ge3\).
\end{lemma}

\begin{proof}
The edges \(vy_1,vy_2\) lie in \(S_\omega\subseteq E(F)\) by (P2). For
\(vx\), take the covering geodesic of (P3) with \(i=1\): a selected
geodesic \(Q=w_0w_1w_2w_3w_4\) of some atom \(ab\in\mathcal A\) with
\(x,y_1\in V(Q)\) and \(v\notin V(Q)\). The vertices \(x,y_1\in
N_B(v)\) lie on the shore opposite \(v\), and the positions of \(Q\)
alternate shores, with \(w_0=a\) and \(w_4=b\) on one shore. If \(a,b\)
lie on the shore of \(v\), then \(x\) and \(y_1\) occupy the two odd
positions \(w_1,w_3\). If \(a,b\) lie opposite \(v\), then \(x,y_1\)
occupy even positions; they cannot be \(\{w_0,w_4\}\), since then
\(xy_1=ab\in M\) and \(vx,vy_1\in B\) would form a triangle. In every
case \(x\) and \(y_1\) occur at distance two along \(Q\), separated by
a middle vertex \(w\), and \(w\ne v\) because \(v\notin V(Q)\).
Replacing \(w\) by \(v\) in \(Q\) yields a walk from \(a\) to \(b\) of
length four whose two new edges \(vx,vy_1\) lie in \(B\) and whose
vertices are distinct; it is therefore a \(B\)-geodesic of \(ab\),
since \(d_B(a,b)=4\). Its edges lie in
\(\supp_B(ab)\subseteq E(F)\) by the definition of the support. Hence
\(vx\in E(F)\). The three edges are distinct because \(x,y_1,y_2\)
are.
\end{proof}

\begin{lemma}[One shore, pairwise nonadjacent]\label{lem:t3-disjoint}
Let \(v_1,\dots,v_k\), \(\omega_1,\dots,\omega_k\) be a balanced rotor
on \((F,\mathcal A)\). Then all rotating owners lie on one shore, any
two of them are nonadjacent in \(G\), and no edge of \(F\) is incident
to two owners. In particular, the sets of \(F\)-edges incident to the
several owners are pairwise disjoint.
\end{lemma}

\begin{proof}
Consider the move (R2) from \(\omega_i\) to \(\omega_{i+1}\), at an
atom \(ab\). Both geodesics \(a\,x\,v_{i+1}\,y\,b\) and
\(a\,x\,v_i\,y\,b\) are five-vertex alternating paths whose endpoints
\(a,b\), being the endpoints of a monochromatic edge, lie on one shore;
the middle position therefore lies on the endpoint shore as well.
Hence \(v_i\) and \(v_{i+1}\) lie on a common shore, and around the
cycle of moves all owners lie on one shore. Consequently
\(v_iv_{i+1}\notin B\); and if \(v_iv_{i+1}\in M\), then, since
\(xv_i,xv_{i+1}\in B\) by the two displayed geodesics, the vertices
\(x,v_i,v_{i+1}\) would form a triangle in \(G\). So consecutive
owners are nonadjacent. For \(k\le3\) every pair of owners is
consecutive in the cyclic order, so no two owners are adjacent, and an
edge of \(G\) incident to two owners would be exactly such an edge.
\end{proof}

\begin{theorem}[First closure: support-edge counting]\label{thm:t3-counting}
A nine-atom support circuit carries no balanced rotor.
\end{theorem}

\begin{proof}
Suppose \((F,\mathcal A)\) with \(|\mathcal A|=9\) carries a balanced
rotor; by Corollary~\ref{cor:graph}(a), \(|E(F)|=8\). Each owner
\(v_i\) has at least three incident edges in \(F\)
(Lemma~\ref{lem:t3-ownerdeg}, applied at \(\omega_i\); note that
\(E(F)\) does not depend on the selection), and the incident sets of
distinct owners are disjoint (Lemma~\ref{lem:t3-disjoint}).

If \(k=3\), the three disjoint triples already give
\(|E(F)|\ge9>8\), a contradiction.

If \(k=2\), write \(v=v_1\), \(m=v_2\). The two incident triples give
six edges of \(F\), and we exhibit three further edges of \(F\)
incident to neither owner. By (P1) at \(v\), the three monochromatic
edges \(vb_1,vb_2,vb_3\) at \(v\) are atoms; the vertices
\(b_1,b_2,b_3\) are distinct, lie on the shore of \(v\), and none
equals \(m\), since \(vm\notin E(G)\) by
Lemma~\ref{lem:t3-disjoint} while \(vb_i\in M\). By
Corollary~\ref{cor:graph}(c), \(d_F(v,b_i)=4\); choose an
\(F\)-geodesic \(P_i\) from \(v\) to \(b_i\) for each \(i\). Two cases.

\emph{Some \(v\)--\(b_i\) geodesic in \(F\) avoids \(m\).} Choose
\(P_i\) to be such a geodesic. Its last three edges are three distinct
edges of \(F\); none is incident to \(v\), which occurs only at the
first position of the path, and none is incident to \(m\), which does
not occur at all.

\emph{Every \(v\)--\(b_i\) geodesic in \(F\) contains \(m\), for each
\(i\).} Fix any choice of \(P_i=v\,p_i\,h_i\,q_i\,b_i\). Since \(m\)
lies on the shore of \(v\), it occupies an even position of \(P_i\);
it is neither position \(0\) (which is \(v\ne m\)) nor position \(4\)
(which is \(b_i\ne m\)), so \(m=h_i\) for every \(i\). The three final
edges \(q_ib_i\) then avoid both owners: \(q_i\), at an odd position,
lies opposite the shore of \(v\) and \(m\), and \(b_i\notin\{v,m\}\).
They are pairwise distinct because each has exactly one endpoint,
namely \(b_i\), on the shore of \(v\), and the \(b_i\) are distinct.

In either case \(E(F)\) contains three edges of \(F\) meeting neither
\(v\) nor \(m\), disjoint from the six incident edges. Hence
\(|E(F)|\ge3+3+3=9>8\), a contradiction.
\end{proof}

\subsection{Second closure: coverage emptiness on the nine-atom
footprint}

The second proof does not count support edges. It shows that on the
unique nine-atom footprint the covered-star clause (P3) can never be
satisfied, at any vertex and any selection. The row-level statement is
general.

\begin{theorem}[Row exclusion at a closed star]\label{thm:t3-emptystar}
Let \((F,\mathcal A)\) be a support circuit and let \(z\) be a vertex
such that every atom of \(\mathcal A\) has both endpoints on the shore
of \(z\), and every atom has at least one endpoint in
\(\{z\}\cup N_M(z)\), where \(N_M(z)\) is the set of monochromatic
neighbours of \(z\). Let \(\omega\) be a selection and let
\(x,s\in N_B(z)\) be distinct with \(zx\notin S_\omega\). Then no
selected geodesic of \(\omega\) contains both \(x\) and \(s\).
\end{theorem}

\begin{proof}
Suppose the selected geodesic \(Q=a\,q_1\,q_2\,q_3\,b\) of an atom
\(ab\in\mathcal A\) contains both \(x\) and \(s\). The endpoints
\(a,b\) lie on the shore of \(z\) by hypothesis, so the even positions
\(a,q_2,b\) of \(Q\) lie on the shore of \(z\), while
\(x,s\in N_B(z)\) lie on the opposite shore. Hence
\(\{x,s\}=\{q_1,q_3\}\). By the second hypothesis applied to \(ab\),
one of the following holds.

\emph{\(z=a\).} If \(x=q_1\), then \(zx=aq_1\in E(Q)\subseteq
S_\omega\), contradicting \(zx\notin S_\omega\). If \(x=q_3\), then
\(zx\in B\) and \(a\,q_3\,b\) is a \(B\)-path of length two,
contradicting \(d_B(a,b)=4\).

\emph{\(z=b\).} Symmetric, exchanging \(q_1\) and \(q_3\).

\emph{\(za\in M\).} The vertex \(q_1\) is one of \(x,s\), so
\(zq_1\in B\); also \(aq_1\in E(Q)\subseteq B\) and \(za\in M\). Then
\(z,a,q_1\) form a triangle in \(G\), a contradiction.

\emph{\(zb\in M\).} Symmetric, with the triangle \(z,b,q_3\).
\end{proof}

\begin{corollary}[Second closure: coverage emptiness]\label{cor:t3-second}
No vertex of a nine-atom support circuit has the rotating owner
profile at any selection. In particular, a nine-atom support circuit
carries no balanced rotor.
\end{corollary}

\begin{proof}
By the exact footprint classification of the companion paper
\cite{Ferudun26supports}, a nine-atom support circuit has, up to
isomorphism, the support graph
\[
 \{l_1,l_2,l_3\}-c_1-c_2-c_3-\{r_1,r_2,r_3\},
 \qquad
 \mathcal A=\{l_ir_j:1\le i,j\le3\},
\]
the double star whose atom set is a \(K_{3,3}\) on the six vertices
\(l_i,r_j\). Every atom is monochromatic, so its endpoints share a
shore, and since the atom graph \(K_{3,3}\) is connected, all six
vertices \(l_i,r_j\) lie on one shore.

Suppose a vertex \(z\) has the rotating owner profile at a selection
\(\omega\). By (P1) the three monochromatic edges at \(z\) are atoms,
so \(z\) is an endpoint of atoms; say \(z=l_i\) (the case \(z=r_j\) is
symmetric). Then every atom \(l_ar_b\) has the endpoint
\(r_b\in N_M(z)\), because \(l_ir_b\in\mathcal A\subseteq M\); and
every atom has both endpoints on the shore of \(z\). Thus
Theorem~\ref{thm:t3-emptystar} applies to \(z\): since
\(zx\notin S_\omega\) by (P2), no selected geodesic contains both
\(x\) and \(y_1\). This contradicts (P3).
\end{proof}

\begin{remark}[Independence of the two closures]\label{rem:t3-independent}
Theorem~\ref{thm:t3-counting} uses no classification and applies to
every nine-atom circuit shape, whatever its footprint; it does not use
coverage emptiness. Corollary~\ref{cor:t3-second} counts nothing, but
consumes the footprint classification of \cite{Ferudun26supports},
which is computer-assisted with two independent exact verifiers.
Either result alone excludes balanced rotors on nine atoms; the second
is in one respect stronger, excluding even a single profiled vertex
rather than only a cyclic family.
\end{remark}

\begin{remark}[Formalization]\label{rem:t3-lean}
Both closures are verified in Lean~4 in the development
\texttt{Erdos23Delta0} accompanying this project, in the following
division of labour. All named declarations compile with axioms exactly
\texttt{propext}, \texttt{Classical.choice}, \texttt{Quot.sound}; none
uses \texttt{sorry} or \texttt{native\_decide}.
\begin{enumerate}[label=\textup{(\roman*)}]
\item For Theorem~\ref{thm:t3-counting}, the finite counting core is
the module \texttt{Erdos23Delta0.Gamma.R43SupportIncidence} (source
SHA-256 prefix \texttt{66649670}). Abstract carriers package an
eight-edge support set together with, per owner, a three-element
incident subset with explicit distinctness witnesses, the pairwise
disjointness facts, and, for two owners, the external triple; the
kernel-checked theorems
\texttt{fullyCoveredLiveStar\_fullSupportDegree\_ge\_three},
\texttt{twoRotatingOwners\_force\_nine\_supportEdges},
\texttt{threeRotatingOwners\_force\_nine\_supportEdges}, and
\texttt{no\_t3\_balancedDeficiencyRotor} derive the contradiction
\(9\le|E(F)|=8\). The graph-level inputs, namely
Lemmas~\ref{lem:t3-ownerdeg} and~\ref{lem:t3-disjoint} and the
two-case argument in the proof of Theorem~\ref{thm:t3-counting}, are
proved on paper above and enter the formal statement as hypotheses of
the carriers.
\item For Theorem~\ref{thm:t3-emptystar}, the formalization is at
graph level, over the development's literal edge-list graphs, explicit
five-vertex checked rows, and executable pair counts. The row-local
statement is \texttt{bad\_star\_cover\_row\_impossible} in
\texttt{Erdos23Delta0.Gamma.BadStarCoverFreeness} (SHA-256 prefix
\texttt{afd944ea}); its endpoint cases use the checked-row inducedness
lemma \texttt{checkedRow\_blue\_cooccur\_implies\_pathEdge} of
\texttt{Erdos23Delta0.Gamma.SaturatedRotorSupportPersistence}. The
selection-level statement, with the closed-star hypotheses of
Theorem~\ref{thm:t3-emptystar} as a structure and the conclusion that
the production pair count of \(x,s\) vanishes, is
\texttt{pairCount\_eq\_zero\_of\_closedBadStar} in
\texttt{Erdos23Delta0.Gamma.K33BadStarPairCountZero} (SHA-256 prefix
\texttt{35456523}).
\item The disjoint-star count of Remark~\ref{rem:t3-fourowners} has
its finite core in
\texttt{Erdos23Delta0.Gamma.CutTightActiveRotorIncidence} (SHA-256
prefix \texttt{0814a665}): \texttt{twelve\_le\_ambient\_card} and
\texttt{not\_ambient\_card\_nine}.
\end{enumerate}
\end{remark}

\begin{remark}[Scope]\label{rem:t3-scope}
Both theorems close only the nine-atom window. At the next window,
sixteen atoms and \(|E(F)|=15\), the analogous counts give two owners
of support degree at least \(t=4\) plus an external family of size
\(t\), in total \(3t=12\le15\) edges, and no contradiction results;
the exclusion of larger windows requires different arguments and is
not addressed in this section. As with the static obstructions of the
companion paper, these results constrain proof strategies: they yield
no upper bound on \(\bip(G)\) and have no bearing on the truth of the
Erd\H{o}s \(n^2/25\) conjecture.
\end{remark}
